\documentclass[a4paper]{article}
\usepackage{amsfonts}
\usepackage{amsmath}
\usepackage{amssymb}
\usepackage{graphicx}     

\begin{document}
  
\noindent \large Auteur: Anna Voronovska \\
\noindent \large Studierichting: Informatica\\
\noindent \large Oefeningenreeks 4A nr. 47.2\\

\medskip

\normalsize



$\displaystyle\int \dfrac{3x-1}{(x^2+4x+5)^3}dx = \displaystyle\int \dfrac{3x-1}{((x+2)^2 + 1)^3}dx$\\

Stel $t = x+2 \Rightarrow dt = dx \Rightarrow x = t-2$:\\

$\displaystyle\int \dfrac{3(t-2)-1}{(t^2+1)^3}dt = \displaystyle\int \dfrac{3t-7}{(t^2+1)^3}dt$\\

$= 3\displaystyle\int \dfrac{t}{(t^2+1)^3}dt - 7\displaystyle\int \dfrac{1}{(t^2+1)^3}dt$\\

Eerste integraal:\\

$u = t^2+1 \Rightarrow du = 2tdt$:\\

$3\displaystyle\int \dfrac{t}{(t^2+1)^3}dt = \dfrac{3}{2}\displaystyle\int u^{-3}du = \dfrac{3}{2} \cdot \dfrac{u^{-2}}{-2} = -\dfrac{3}{4(t^2+1)^2}$\\

Tweede integraal formule:\\

$I_n = \displaystyle\int \dfrac{dt}{(t^2+1)^n} = \dfrac{t}{2(n-1)(t^2+1)^{n-1}} + \dfrac{2n-3}{2(n-1)} I_{n-1}$\\

Voor $n=2$:\\

$I_2 = \dfrac{t}{2(t^2+1)} + \dfrac{1}{2}\operatorname{Bgtan} t$\\

Voor $n=3$:\\

$I_3 = \dfrac{t}{4(t^2+1)^2} + \dfrac{3}{4}I_2 = \dfrac{t}{4(t^2+1)^2} + \dfrac{3t}{8(t^2+1)} + \dfrac{3}{8}\operatorname{Bgtan}t$\\

Samen:\\

$\displaystyle\int \dfrac{3x-1}{(x^2+4x+5)^3}dx $\\

$= -\dfrac{3}{4(t^2+1)^2} - 7\left(\dfrac{t}{4(t^2+1)^2} + \dfrac{3t}{8(t^2+1)} + \dfrac{3}{8}\operatorname{Bgtan}t\right) + c$\\

$= \dfrac{-3-7t}{4(t^2+1)^2} - \dfrac{21t}{8(t^2+1)} - \dfrac{21}{8}\operatorname{Bgtan} t + c$\\

Substitutie terug naar $x$:\\

$\displaystyle\int \dfrac{3x-1}{(x^2+4x+5)^3}dx$\\

$= \dfrac{-7x-17}{4(x^2+4x+5)^2} - \dfrac{21(x+2)}{8(x^2+4x+5)} - \dfrac{21}{8}\operatorname{Bgtan}(x+2) + c$



\begin{thebibliography}{999}
%\bibitem{a} Referentie
\end{thebibliography}
\end{document} 
